% GENERATED FILE. Edit canonical lessons, then run npm run generate:books.
% canonical-source-hash: fnv1a64:265d3ef9a4120193
% canonical-lessons: PT-C413-abrigo, PT-C413-terminal, PT-C413-camioneta, PT-C413-itinerario, PT-C413-apear

\chapter{Na paragem}
\label{ch:pt-na-paragem}
\hlchaptermodality{413}

\begin{chapteropening}
\textbf{By the end of this chapter:} \emph{I can wait for a bus or a coach: the shelter, the terminus, the coach, the route and getting off.}

Complete the last lesson of chapter 413: I can wait for a bus or a coach: the shelter, the terminus, the coach, the route and getting off.
\end{chapteropening}

\section[o abrigo]{o abrigo — the shelter}

\label{lesson:PT-C413-abrigo}



\begin{quote}
Before the new one: say the Portuguese for to pedal, then the Portuguese for the parking meter.
\end{quote}

\subsection*{You'll want to know: o abrigo}
\textbf{o abrigo} — \textquotedblleft{}the shelter\textquotedblright{}.

\begin{quote}
Esperamos o autocarro no abrigo da paragem.
\end{quote}

\begin{quote}
\emph{We wait for the bus in the bus shelter.}
\end{quote}

\emph{Ficar ao abrigo da chuva} — to stay out of the rain.

\begin{cousinweb}
\emph{abrigo} is built on \emph{abrigar} 'to shelter' (EtymDB). \emph{Abrigar} is usually traced to Latin \emph{aprīcārī} 'to warm oneself in the sun' (Lewis \& Short), but that is not settled.
\end{cousinweb}

\begin{grammarlens}[title={What you've built}]
One more word for telling, explaining and describing at B1.
\end{grammarlens}

\subsection*{Your turn}
\begin{itemize}
\raggedright
  \item \emph{Say it:} \emph{o abrigo}
  \item \emph{Say it:} \emph{o abrigo}, once more
  \item \emph{Say it:} a sentence of your own with \emph{o abrigo}
  \item \emph{Recall:} say the Portuguese for to pedal, then the Portuguese for the parking meter, then say \emph{o abrigo} again
\end{itemize}

\begin{tcolorbox}[breakable,colback=teal!4,colframe=teal!35!black,title={Before you move on}]
What does o abrigo mean? (\textquotedblleft{}The shelter\textquotedblright{}.) Say it once more.
\end{tcolorbox}
\section[o terminal]{o terminal — the terminus, the terminal}

\label{lesson:PT-C413-terminal}



\begin{quote}
Before the new one: say the Portuguese for the parking meter, then the Portuguese for the shelter.
\end{quote}

\subsection*{You'll want to know: o terminal}
\textbf{o terminal} — \textquotedblleft{}the terminus, the terminal\textquotedblright{}.

\begin{quote}
O autocarro acaba a viagem no terminal.
\end{quote}

\begin{quote}
\emph{The bus ends its journey at the terminus.}
\end{quote}

Also an airport terminal: \emph{o terminal 2}.

\begin{cousinweb}
\emph{terminal} is a learned borrowing from Late Latin \emph{terminālis} 'of a boundary', from \emph{terminus} 'boundary, end' (EtymDB).
\end{cousinweb}

\begin{grammarlens}[title={What you've built}]
One more word for telling, explaining and describing at B1.
\end{grammarlens}

\subsection*{Your turn}
\begin{itemize}
\raggedright
  \item \emph{Say it:} \emph{o terminal}
  \item \emph{Say it:} \emph{o terminal}, once more
  \item \emph{Say it:} a sentence of your own with \emph{o terminal}
  \item \emph{Recall:} say the Portuguese for the parking meter, then the Portuguese for the shelter, then say \emph{o terminal} again
\end{itemize}

\begin{tcolorbox}[breakable,colback=teal!4,colframe=teal!35!black,title={Before you move on}]
What does o terminal mean? (\textquotedblleft{}The terminus, the terminal\textquotedblright{}.) Say it once more.
\end{tcolorbox}
\section[a camioneta]{a camioneta — the coach, the long-distance bus}

\label{lesson:PT-C413-camioneta}



\begin{quote}
Before the new one: say the Portuguese for the shelter, then the Portuguese for the terminus.
\end{quote}

\subsection*{You'll want to know: a camioneta}
\textbf{a camioneta} — \textquotedblleft{}the coach, the long-distance bus\textquotedblright{}.

\begin{quote}
Fomos para o sul de camioneta.
\end{quote}

\begin{quote}
\emph{We went south by coach.}
\end{quote}

In Portugal \emph{a camioneta} runs between towns and \emph{o autocarro} inside a town. In Brazil: \emph{o ônibus}.

\begin{cousinweb}
\emph{camioneta} is a small form of \emph{camião} 'lorry', which was borrowed from French \emph{camion} 'truck' (EtymDB); the French word's origin is not settled.
\end{cousinweb}

\begin{grammarlens}[title={What you've built}]
One more word for telling, explaining and describing at B1.
\end{grammarlens}

\subsection*{Your turn}
\begin{itemize}
\raggedright
  \item \emph{Say it:} \emph{a camioneta}
  \item \emph{Say it:} \emph{a camioneta}, once more
  \item \emph{Say it:} a sentence of your own with \emph{a camioneta}
  \item \emph{Recall:} say the Portuguese for the shelter, then the Portuguese for the terminus, then say \emph{a camioneta} again
\end{itemize}

\begin{tcolorbox}[breakable,colback=teal!4,colframe=teal!35!black,title={Before you move on}]
What does a camioneta mean? (\textquotedblleft{}The coach, the long-distance bus\textquotedblright{}.) Say it once more.
\end{tcolorbox}
\section[o itinerário]{o itinerário — the route, the itinerary}

\label{lesson:PT-C413-itinerario}



\begin{quote}
Before the new one: say the Portuguese for the terminus, then the Portuguese for the coach.
\end{quote}

\subsection*{You'll want to know: o itinerário}
\textbf{o itinerário} — \textquotedblleft{}the route, the itinerary\textquotedblright{}.

\begin{quote}
O itinerário da carreira passa pelo hospital.
\end{quote}

\begin{quote}
\emph{The route of the bus line goes past the hospital.}
\end{quote}

\emph{O itinerário} of a trip is also its plan, day by day.

\begin{cousinweb}
\emph{itinerário} is a learned borrowing from Late Latin \emph{itinerārium} 'travel guide', built on \emph{iter, itineris} 'journey' (Lewis \& Short).
\end{cousinweb}

\begin{grammarlens}[title={What you've built}]
One more word for telling, explaining and describing at B1.
\end{grammarlens}

\subsection*{Your turn}
\begin{itemize}
\raggedright
  \item \emph{Say it:} \emph{o itinerário}
  \item \emph{Say it:} \emph{o itinerário}, once more
  \item \emph{Say it:} a sentence of your own with \emph{o itinerário}
  \item \emph{Recall:} say the Portuguese for the terminus, then the Portuguese for the coach, then say \emph{o itinerário} again
\end{itemize}

\begin{tcolorbox}[breakable,colback=teal!4,colframe=teal!35!black,title={Before you move on}]
What does o itinerário mean? (\textquotedblleft{}The route, the itinerary\textquotedblright{}.) Say it once more.
\end{tcolorbox}
\section[apear]{apear — to get off (a vehicle); to dismount}

\label{lesson:PT-C413-apear}



\begin{quote}
Before the new one: say the Portuguese for the coach, then the Portuguese for the route.
\end{quote}

\subsection*{You'll want to know: apear}
\textbf{apear} — \textquotedblleft{}to get off (a vehicle); to dismount\textquotedblright{}.

\begin{quote}
Apeei-me na paragem do mercado.
\end{quote}

\begin{quote}
\emph{I got off at the market stop.}
\end{quote}

Mostly reflexive and a bit formal: \emph{apear-se do comboio}. Everyday speech also says \emph{sair do autocarro}.

\begin{cousinweb}
\emph{apear} is built from \emph{a-}, \emph{pé} 'foot' and \emph{-ar} (EtymDB): to put your feet on the ground. \emph{Pé} is inherited from Latin \emph{pedem}.
\end{cousinweb}

\begin{grammarlens}[title={What you've built}]
One more word for telling, explaining and describing at B1.
\end{grammarlens}

\subsection*{Your turn}
\begin{itemize}
\raggedright
  \item \emph{Say it:} \emph{apear}
  \item \emph{Say it:} \emph{apear}, once more
  \item \emph{Say it:} a sentence of your own with \emph{apear}
  \item \emph{Recall:} say the Portuguese for the coach, then the Portuguese for the route, then say \emph{apear} again
\end{itemize}

\begin{tcolorbox}[breakable,colback=teal!4,colframe=teal!35!black,title={Before you move on}]
What does apear mean? (\textquotedblleft{}To get off (a vehicle); to dismount\textquotedblright{}.) Say it once more.
\end{tcolorbox}
